% Regression: p3_blocking — imported current-behavior fixture from the handoff corpus.
% p3_blocking — arXiv:1804.04964 (normal PEPS FT), proof of Theorem 3, p. 15.
% Region-to-MPS reduction: block the 7x7 TN into three injective parts
% around one edge.
%
% Formula: choose the horizontal edge e = <(4,3),(4,4)>; region A_1 =
% rows 2-4 x cols 2-3 (a 3x2 injective block containing the west endpoint),
% A_2 = rows 4-5 x cols 4-6 (a 2x3 injective block containing the east
% endpoint), A_3 = the rest.  Contracting each region gives the closed
% (trace-boundary) three-partite injective MPS A_1 A_2 A_3, with the
% distinguished edge surviving as the bond A_1 - A_2, so the injective-MPS
% edge-isomorphism lemma applies to e.
%
% Hue mapping noted below: the paper keys A_1/A_2/A_3 red/blue/black and
% the edge green; tenkz's semantic slots give selected=red, secondary=blue,
% distinguished edge=Action red.  The paper's third hue (green edge) and
% the chain's per-cell hues have no tenkz primitive.
\documentclass[varwidth=24cm, border=6pt]{standalone}
\usepackage{amsmath, amssymb}
\usepackage{tenkz}
\begin{document}
$
\begin{tenkzlattice}[rows=7, cols=7, boundary legs]
  % A_1: outline-only 3x2 block around the west endpoint (4,3)
  \tnregion[slot=selected, outline]{(2-4,2-3)}
  % A_2: outline-only 2x3 block around the east endpoint (4,4)
  \tnregion[slot=secondary, outline]{(4-5,4-6)}
  % the distinguished edge e between the two endpoints
  \tnedge{(4,3)-(4,4)}
\end{tenkzlattice}
\;\Rightarrow\;
% the reduced chain: closed three-partite MPS (trace boundary), physical
% legs up; the paper draws the closure as the rectangle under the chain
\begin{tenkz}[physical=up, periodic]
  \tn{A_1} & \tn{A_2} & \tn{A_3}
\end{tenkz}
$
\end{document}
